% GENERATED_BY: run_oc_core_monograph_translation_rework_dispatch_v1.py
% AI_ASSISTED_TRANSLATION_DRAFT: TRUE
% THEOREM_REVIEW_STATUS: APPROVED
% THEOREM_REVIEW_MODE: FORMAL_AI_GATE__CODEX_CLI_CHATGPT
% THEOREM_REVIEW_REVIEWER_ID: CODEX_CLI_CHATGPT__AUTO_DISPATCH_20260406
% THEOREM_REVIEW_AUTHORITY_CLASS: FINAL_ADVERSARIAL_EXTERNAL_LLM_REVIEWER
% THEOREM_REVIEWED_AT: 2026-04-14T14:17:20Z
% SOURCE_LANGUAGE: EN
% TARGET_LANGUAGE: RU
% SOURCE_EN_REF: content/02_background.tex
% FILE: content/02_background_ru.tex

\section{Предпосылки и мотивация}
\label{sec:background}

В этой главе вводятся структурные и математические предпосылки, необходимые для Онтологии континуумов (OC).
Ее цель состоит не в историческом изложении, а в представлении минимального понятийного аппарата, лежащего в основе единой рамки.
Все ключевые понятия, а именно континуумы, оси, пороги, границы, потенциалы, потоки, циклы и континуумность, вводятся в форме, не зависящей от предметной области, независимо от того, является ли целевая система физической, химической, биологической, когнитивной, социальной или метатеоретической.

\subsection{Континуум как структурный объект}

В OC \emph{континуум} определяется не пространственной гладкостью и не физической протяженностью, а набором структурных условий.
Система квалифицируется как континуум \(K\) тогда и только тогда, когда она обладает:

\begin{itemize}
    \item непустым допустимым пространством состояний \(\Omega(K)\);
    \item конечным множеством осей \(A(K)=\{A_1,\dots,A_n\}\), каждая из которых представляет независимое и несовместимое структурное различие;
    \item потенциалами \(P(t)\), фиксирующими энергетические, химические, информационные, биологические, когнитивные или институциональные ограничения;
    \item потоками \(J(t)\), которые преобразуют эти потенциалы и направляют эволюцию;
    \item ландшафтом порогов \(\Theta(K)\), разделяющим качественно различные динамические режимы;
    \item границей \(\partial\Omega(K)\), на которой пороги достигают насыщения;
    \item репертуаром устойчивых циклов \(C(K)\), поддерживающих организацию и удерживающих динамику вдали от \(\partial\Omega(K)\);
    \item мерой континуумности \(k(K,t)\), количественно выражающей жизнеспособность и структурную целостность;
    \item метапространством вложения \(M\), для которого \(\Omega(K)\subseteq\Omega(M)\) и \(A(K)\subseteq A(M)\).
\end{itemize}

Это определение намеренно избегает предметно-специфических деталей.
Физические поля, сети химических реакций, протоклетки, нейронные ансамбли, системы представлений, социальные институты и научные теории все считаются континуумами, если они удовлетворяют этим структурным критериям.

\subsection{Оси и размерность}

Оси \(A(K)\) определяют эффективную размерность континуума.
Каждая ось представляет \emph{несовместимое различие}, то есть различение, которое не может быть реконструировано из уже имеющихся осей.
Примеры включают энергетические оси, градиентные оси, мембранные оси, оси возбуждения, когнитивные оси, оси социальных различий и выразительные оси более высокого уровня.

Размерность подчиняется \emph{принципу монотонности} OC:
\[
\dim(K_{x+1}) > \dim(K_x)
\quad\text{всякий раз, когда возникает новый континуум}.
\]

Новая размерность появляется только тогда, когда:

\begin{enumerate}
    \item возникает новый класс различий, который не может быть выражен в пределах \(\mathrm{span}(A(K_x))\);
    \item структурное напряжение \(T(K_x)\) превышает порог размерности \(\Theta_{\mathrm{dim}}(K_x)\);
    \item пространство вложения \(M_x\) предоставляет ось \(A_{\mathrm{new}}\in A(M_x)\setminus A(K_x)\).
\end{enumerate}

Следовательно, размерность является не свободным параметром, а структурным инвариантом, навязываемым несовместимостью и наличием подходящих осей в окружающем метапространстве.

\subsection{Границы и пороги}

Континуум ограничен своим ландшафтом порогов \(\Theta(K)\).
Пороги являются функциями
\[
\Theta_k : \Omega(K) \to \mathbb{R},
\qquad
\Theta_k(s)\le 0 \text{ для } s\in\Omega(K),
\]
и подразделяются на структурные классы:

\begin{itemize}
    \item \textbf{Пороги существования} \(\Theta_{\mathrm{exist}}\): условия, при которых \(\Omega(K)\neq\emptyset\);
    \item \textbf{Пороги устойчивости} \(\Theta_{\mathrm{stab}}\): обеспечивают нерасходящиеся потоки;
    \item \textbf{Критические пороги} \(\Theta_{\mathrm{crit}}\): отмечают качественные переходы;
    \item \textbf{Пороги размерности} \(\Theta_{\mathrm{dim}}\): управляют возникновением новых осей;
    \item \textbf{Пороги смерти} \(\Theta_{\mathrm{death}}\): за их пределами не остается допустимых состояний.
\end{itemize}

Структурная граница определяется как
\[
\partial\Omega(K) = \{\, s\in\Omega(K) \mid \exists k: \Theta_k(s)=0 \,\}.
\]

Границы кодируют полную структуру пределов континуума.
Их динамика описывается оператором
\[
\frac{d}{dt}\,\partial\Omega = R(\partial\Omega, P, J, \Theta),
\]
который описывает расширение, сжатие или бифуркацию \(\partial\Omega\) при изменении потенциалов, потоков или порогов.
Рождение, сохранение и коллапс соответствуют различным режимам \(R\).

\subsection{Потенциалы и потоки}

Потенциалы \(P(t)\) кодируют внутренние ограничения и движущие силы.
К характерным примерам относятся:

\begin{itemize}
    \item энергетические ландшафты, поля и параметры порядка (физика);
    \item концентрации, уровни pH, редокс-потенциалы (химия);
    \item мембранные и электрохимические потенциалы (биология);
    \item предиктивные и репрезентационные потенциалы (когниция);
    \item нормативные и институциональные потенциалы (социальные системы).
\end{itemize}

Потоки описывают эволюцию потенциалов:
\[
\frac{dP}{dt} = J(t).
\]

В OC различаются:

\begin{itemize}
    \item \textbf{Поддерживающие потоки} \(J_{\mathrm{support}}\): поддерживают циклы и стабилизируют \(k(t)\);
    \item \textbf{Критические потоки} \(J_{\mathrm{critical}}\): подталкивают систему к \(\Theta_{\mathrm{crit}}\) или \(\Theta_{\mathrm{dim}}\), либо через них;
    \item \textbf{Деструктивные потоки} \(J_{\mathrm{kill}}\): разрушают циклы или вынуждают систему пересечь \(\Theta_{\mathrm{death}}\).
\end{itemize}

Эти потоки входят в универсальные операторы эволюции \(F, G, H, Q, R, S, U\), которые задают, как совместно эволюционируют оси, потенциалы, пороги, потоки, циклы и границы.

\subsection{Континуумность}

Мера \(k(K,t)\) определяет, функционирует ли система как \emph{живой континуум}.
В Core~1.3 принимается единое определение, закодированное оператором \(U\):

\[
k(K,t) =
\chi_{\Omega}(K,t)\;
S_{\mathrm{axes}}(K,t)\;
S_{\mathrm{cycles}}(K,t)\;
S_{\mathrm{flows}}(K,t)\;
S_{\mathrm{coh}}(K,t),
\]

где:

\begin{itemize}
    \item \(\chi_{\Omega}(K,t)=1\), если \(\Omega(K)\neq\emptyset\), и \(0\) в противном случае;
    \item \(S_{\mathrm{axes}}\) измеряет выразительную адекватность осей;
    \item \(S_{\mathrm{cycles}}\) измеряет силу стабилизирующих циклов;
    \item \(S_{\mathrm{flows}}\) отражает баланс поддерживающих и деструктивных потоков;
    \item \(S_{\mathrm{coh}}\) кодирует глобальную структурную когерентность.
\end{itemize}

Континуум является живым в точности тогда, когда
\[
k(K,t) > 0.
\]
Смерть соответствует \(k(K,t)\to 0\) и \(\Omega(K)\to\emptyset\), независимо от пути.

\subsection{Мотивация для Core}

Core~1.3 объединяет все структурные компоненты OC в вертикально согласованном опорном изложении.
Его цели состоят в следующем:

\begin{itemize}
    \item установить каноническую структуру файлов и репозитория;
    \item консолидировать аксиоматику \(K_0\) и построение \(K_1\);
    \item унифицировать трактовку границ и порогов;
    \item формализовать потенциалы, потоки, циклы и операторы эволюции на согласованном языке;
    \item интегрировать полную иерархию \(K_0\)–\(K_{10}\) вместе с метапространствами вложения \(M_x\).
\end{itemize}

\subsection{Историческая мотивация}

OC возникла из попыток объяснить, почему разнообразные системы, такие как фазовые переходы, каталитическое замыкание, стабилизация протоклеток, нейронное возбуждение, институциональная когерентность и цивилизационный коллапс, разделяют одни и те же организационные инварианты.
Во всех предметных областях вновь и вновь проявляются четыре структуры:

\begin{enumerate}
    \item пороги, управляющие допустимостью и устойчивостью;
    \item возникновение новых размерностей при структурном напряжении;
    \item циклы, стабилизирующие потоки;
    \item коллапс при нарушении порогов и исчезновении \(\Omega(K)\).
\end{enumerate}

OC формализует эти инварианты в рамках единой структурной онтологии.

\subsection{Объем будущего материала о предпосылках}

В будущие расширения этой главы войдут:

\begin{itemize}
    \item математические предварительные сведения о потенциалах, потоках и устойчивости циклов;
    \item структурные мотивации для аксиоматики OC;
    \item уточненное описание пространств вложения \(M_x\) и монотонного расширения;
    \item формальные трактовки монотонности, когерентности и устойчивости на разных уровнях;
    \item выводы универсальных операторов \(F,G,H,Q,R,S,U\) из более примитивных предпосылок.
\end{itemize}

В Core~1.3 включены только минимальные предпосылки, необходимые для формальной модели в разделе~\ref{sec:model}.
